\subsection{定义与一般判别法}

定义 1. 设 X 是拓扑空间，Y 是一致空间。称 F(X; Y) 的子集 H 在一点 \(x_{0} \in X\) 处等度连续，如果对 Y 的每个围域 V，存在 \(x_{0}\) 在 X 中的邻域 U，使得 \((f(x_{0}),f(x))\in\mathbf{V}\) 对一切 \(x\in\mathbf{U}\) 及一切 \(f\in\mathbf{H}\) 成立。称 H 是等度连续的，如果它在 X 的每一点都等度连续。

定义 2. 设 X 与 Y 是两个一致空间。称 \(\mathcal{F}(X;Y)\) 的子集 H 是一致等度连续的，如果对 Y 的每个围域 V，存在 X 的围域 U，使得我们就有 \((f(x), f(x')) \in V\) ，只要 \((x, x') \in U\) 且 \(f \in H\) 。

称 X 到 Y 中的映射的族 \((f_{t})_{t\in I}\) 在一点 \(x_0\) 处等度连续（相应地，等度连续、一致等度连续），如果诸 \(f_{t}\) 所成之集在 \(x_0\) 处等度连续（相应地，等度连续、一致等度连续）。

显然，若 \(\mathbf{H}\subset\mathcal{F}(\mathbf{X};\mathbf{Y})\) 在 \(x_{0}\) 处等度连续，则每个 \(f\in H\) 在 \(x_{0}\) 处连续；若 H 等度连续，则每个 \(f\in H\) 在 X 上连续，即 \(\mathbf{H}\subset\mathcal{C}(\mathbf{X};\mathbf{Y})\) 。同样，若 H 一致等度连续（X 是一致空间），则每个 \(f\in H\) 在 X 上一致连续。同样显然的是，一致等度连续的映射集是等度连续的；但一个由一致连续映射组成的集可以是等度连续的而非一致等度连续的（见习题 1；命题 1 的推论 2；以及第 2 号的命题 4）。

例。1) 设 X 是拓扑空间（相应地，一致空间），Y 是一致空间。由 X 到 Y 中的连续（相应地，一致连续）映射所成的每个有限集都是等度连续（相应地，一致等度连续）的。

\begin{enumerate}
\def\labelenumi{\arabic{enumi})}
\setcounter{enumi}{1}
\tightlist
\item
  设 X、Y 是两个度量空间，d（相应地，d'）是 X（相应地，Y）上的度量，并设 k、\(\alpha\) 是两个实数 \textgreater{} 0。则使
\end{enumerate}

\[
d ^ {\prime} (f (x), f (x ^ {\prime})) \leqslant k (d (x, x ^ {\prime})) ^ {\alpha}
\]

对每一对点 \(x, x'\) （属于 \(\mathbf{X}\) ）都成立的全体映射 f: X → Y 所成之集是一致等度连续的。例如，\(\mathbf{X}\) 到 \(\mathbf{Y}\) 的一个子集上的全体等距映射（第 IX 章，§ 2，第 2 号）是一致等度连续的。

* 设 H 是定义在区间 I ⊂ R 上的实值函数所成之集，这些函数在 I 上可微，且满足 \(|f'(x)| \leqslant k\) 对所有 \(x \in I\) 及所有 \(f \in H\) 成立。则 H 是一致等度连续的，因为若 \(x_{1}, x_{2}\) 是 I 的任意两点，则由中值定理，有 \(|f(x_{1}) - f(x_{2})| \leqslant k |x_{1} - x_{2}|\) 对每个 \(f \in H\) 成立。 *

\begin{enumerate}
\def\labelenumi{\arabic{enumi})}
\setcounter{enumi}{2}
\tightlist
\item
  设 G 是拓扑群，Y 是一致空间，\(f: \mathbf{G} \to \mathbf{Y}\) 是一致连续映射{[}G 赋有它的左一致结构（第 III 章，§ 3，第 1 号）{]}。对每个 \(s \in \mathbf{G}\) ，用 \(f_s\) 表示 G 到 Y 中的映射 \(x \to f(sx)\) 。则诸映射 \(f_s(s \in \mathbf{G})\) 所成之集是一致等度连续的，因为关系 \(x^{-1}x' \in V\) 等价于 \((sx)^{-1}(sx') \in V\) 。
\end{enumerate}

命题 1. 设 T 是一个集合，S 是 T 的子集所成之集，Y 是一致空间，X 是拓扑（相应地，一致）空间，f 是 \(\mathbf{T} \times \mathbf{X}\) 到 \(\mathbf{Y}\) 中的一个映射。对每个 \(\mathbf{A} \in \mathfrak{S}\) ，用 \(\mathbf{H}_{\mathbf{A}} \subset \mathcal{F}(\mathbf{X}; \mathbf{Y})\) 表示诸映射 \(x \to f(t, x)\) （ \(t\) 取遍 \(\mathbf{A}\) ）所成之集。则映射 \(x \to f(., x)\) （由 \(\mathbf{X}\) 到 \(\mathcal{F}_{\mathfrak{S}}(\mathbf{T}; \mathbf{Y})\) 中）在一点 \(x_0 \in \mathbf{X}\) 处连续（相应地，一致连续）当且仅当集 \(\mathbf{H}_{\mathbf{A}}\) 在 \(x_0\) 处等度连续（相应地，一致等度连续）对一切 \(\mathbf{A} \in \mathfrak{S}\) 成立。

先考虑 \(\mathfrak{S} = \{\mathrm{T}\}\) ，即 \(\mathcal{F}_{\mathfrak{S}}(\mathrm{T};\mathrm{Y}) = \mathcal{F}_u(\mathrm{T};\mathrm{Y})\) 这一特殊情形。对 Y 的每个围域 V，条件 \((f(.,x),f(.,x'))\in W(\mathrm{V})\) 的意思是 \((f(t,x),f(t,x'))\in \mathrm{V}\) 对一切 \(t\in \mathrm{T}\) 成立。因此，说 \(x\to f(.,x)\) 在 \(x_0\) 处连续（相应地，一致连续）等价于说：对 Y 的每个围域 V，存在 \(x_0\) 在 X 中的邻域 U（相应地，X 的围域 M），使得关系 \(x\in \mathrm{U}\) {[}相应地 \((x,x')\in \mathrm{M}]\) 蕴涵 \((f(t,x),f(t,x_0))\in \mathrm{V}\) {[}相应地 \((f(t,x),f(t,x'))\in \mathrm{V}]\) 对一切 \(t\in \mathrm{T}\) 成立，于是命题由定义 1 与定义 2 得出。在一般情形，依 § 1，第 2 号，我们要表述的是：对每个 \(\mathrm{A}\in \mathfrak{S}\) ，映射 \(x\rightarrow f(.,x)|\mathrm{A}\) （由 X 到 \(\mathcal{F}_u(\mathrm{A};\mathrm{Y})\) 中）在 \(x_0\) 处连续（相应地，一致连续）；由前所述，这等价于说，对每个 \(\mathrm{A}\in \mathfrak{S}\) ，\(\mathrm{H}_{\mathrm{A}}\) 在 \(x_0\) 处等度连续（相应地，一致等度连续）。

把命题 1 应用于 \(\mathbf{T} = \mathbf{H}\) 且 \(f\) 是映射 \((h, x) \to h(x)\) （由 \(\mathbf{H} \times \mathbf{X}\) 到 \(\mathbf{Y}\) 中）的情形，便可把定义 1 与定义 2 改写成有时有用的形式；由于 \(f(., x)\) 是映射 \(h \to h(x)\) （由 \(\mathbf{H}\) 到 \(\mathbf{Y}\) 中），我们得到：

推论 1. 设 X 是拓扑（相应地，一致）空间，Y 是一致空间，H 是 \(\mathcal{F}\) (X; Y) 的子集。对每个 \(x \in \mathbf{X}\) ，用 \(\tilde{x}\) 表示 H 到 Y 中的映射 \(h \to h(x)\) 。则 H 在 \(x_0\) 处等度连续（相应地，一致等度连续）当且仅当映射 \(x \to \tilde{x}\) （由 X 到一致空间 \(\mathcal{F}_u\) (H; Y) 中）在 \(x_0\) 处连续（相应地，一致连续）。

特别地，若 X 是紧的，则 X 到 \(\mathcal{F}_{u}(H;Y)\) 中的每个连续映射都一致连续（第 II 章，§ 4，第 1 号，定理 2）。因此：

推论 2. 设 X 是紧空间，Y 是一致空间。则 \(\mathcal{F}(X;Y)\) 的每个等度连续子集都是一致等度连续的。

现在设有一个集合 T、一个拓扑空间 X、一个一致空间 Y 以及一个映射 \(f: T \times X \to Y\) 。用 \(\tilde{f}\) 表示映射 \(x \to f(., x)\) （由 X 到 \(\mathcal{F}_{u}(T; Y)\) 中），并考虑典型映射 \(\theta: (t, g) \to g(t)\) （由 \(T \times \mathcal{F}_{u}(T; Y)\) 到 Y 中）。显然，图表

\[
\mathrm{T} \times \mathrm{X} \xrightarrow {f} \mathrm{Y}
\]

\[
\iota_ {\mathrm{T}} \times \tilde {f} \searrow \nearrow \theta
\]

\[
\mathbf {T} \times \mathcal {F} _ {u} (\mathbf {T}; \mathbf {Y})
\]

（其中 \(\iota_{\mathbf{T}}\) 是 T 的恒等映射）是交换的。现在设 T 赋有一个拓扑，并且对每个 \(x \in \mathbf{X}\) ，映射 \(f(., x): t \to f(t, x)\) 连续；这时在上图中就可把 \(\mathcal{F}_{u}(\mathrm{T}; \mathrm{Y})\) 换成 \(\mathcal{C}_{u}(\mathrm{T}; \mathrm{Y})\) 。但由 § 1，第 6 号，命题 9 知 \(\theta\) 连续；因此若 \(\tilde{f}\) 连续，则 \(f\) 连续。由于 \(\tilde{f}\) 的连续性可以借助命题 1 来表达，我们得到下列结果：

推论 3. 设 T、X 是拓扑空间，Y 是一致空间，f 是 T × X 到 Y 中的一个映射。则 f 连续，只要下列条件满足：

\begin{enumerate}
\def\labelenumi{\arabic{enumi})}
\item
  对每个 \(x \in \mathbf{X}\) ，偏映射 \(t \to f(t, x)\) 连续。
\item
  当 \(t\) 取遍 \(\mathbf{T}\) 时，诸偏映射 \(x \to f(t, x)\) 组成 \(\mathfrak{F}(\mathbf{X}; \mathbf{Y})\) 的一个等度连续子集。
\end{enumerate}

特别地，取 T 为 \(\mathcal{F}(\mathbf{X};\mathbf{Y})\) 的一个子集 H，取 \(f\) 为典型映射 \((h,x)\to h(x)\) （由 \(\mathbf{H}\times \mathbf{X}\) 到 Y 中）；推论 3 的条件 1）意味着 H 赋有一个比逐点收敛拓扑更细的拓扑，而条件 2）意味着 H 等度连续。因此：

推论 4. 设 X 是拓扑空间，Y 是一致空间，H 是 X 到 Y 中的等度连续映射集。若 H 赋有逐点收敛拓扑，则映射 \((h, x) \to h(x)\) （由 \(H \times X\) 到 Y 中）连续。

更直观地，这表达了下列事实：若 \(h \in \mathbf{H}\) 逐点收敛于 \(h_0 \in \mathbf{H}\) ，且 \(x \in \mathbf{X}\) 收敛于 \(x_0\) ，则 \(h(x)\) 收敛于 \(h_0(x_0)\) 。

推论 5. 设 X 是拓扑空间，Y、Z 是两个一致空间，H 是 Y 到 Z 中的等度连续映射集。若 H、C(X; Y) 与 C(X; Z) 都赋有逐点收敛拓扑，则映射 \((u, v) \to u \circ v\) （由 H × C(X; Y) 到 C(X; Z) 中）连续。

我们必须证明，对每个 \(x \in \mathbf{X}\) ，映射 \((u, v) \to u(v(x))\) （由 \(\mathbf{H} \times \mathcal{C}(\mathbf{X}; \mathbf{Y})\) 到 \(\mathbf{Z}\) 中）连续。而 \(v \to v(x)\) 在 \(\mathbf{H}\) 上连续（ \(\S 1\) ，第 2 号，注 6），且由推论 4 知 \((u, y) \to u(y)\) 是 \(\mathbf{H} \times \mathbf{Y}\) 到 \(\mathbf{Z}\) 中的连续映射；由于 \((u, v) \to u(v(x))\) 是 \((u, y) \to u(y)\) 与 \((u, v) \to (u, v(x))\) 的复合，结论得证。

下面的命题及其推论是命题 1 的推论 3 与推论 4 关于一致等度连续映射集的类似结果：

命题 2. 设 T、X、Y 是一致空间，f 是 T × X 到 Y 中的一个映射。则 f 一致连续当且仅当下列两个条件满足：

\begin{enumerate}
\def\labelenumi{\arabic{enumi})}
\item
  诸映射 \(x \to f(t, x)\) （ \(t \in \mathbf{T}\) ）组成 \(\mathcal{F}(\mathbf{X}; \mathbf{Y})\) 的一个一致等度连续子集。
\item
  诸映射 \(t \to f(t, x)\) （ \(x \in \mathbf{X}\) ）组成 \(\mathcal{F}(\mathbf{T}; \mathbf{Y})\) 的一个一致等度连续子集。
\end{enumerate}

条件的必要性是容易看出的。反之，设它们成立。设 W 是 Y 的一个围域；则存在 T 的围域 U 与 X 的围域 V，使得：1） \((t', t'') \in \mathbf{U}\) 蕴涵：对每个 \(x \in X\) ，

\[
\left(f \left(t ^ {\prime}, x\right), f \left(t ^ {\prime \prime}, x\right)\right) \in W.
\]

\begin{enumerate}
\def\labelenumi{\arabic{enumi})}
\setcounter{enumi}{1}
\tightlist
\item
  \((x', x'') \in \mathbf{V}\) 蕴涵：对每个 \(t \in \mathbf{T}\) ，
\end{enumerate}

\[
(f (t, x ^ {\prime}), f (t, x ^ {\prime \prime})) \in W.
\]

现在显然，关系“ \((t', t'') \in \mathbf{U}\) 且 \((x', x'') \in \mathbf{V}\) ”蕴涵 \((f(t', x'), f(t'', x'')) \in \overset{2}{\mathbf{W}}\) ，由此即得结论。

特别地，取 T 为 \(\mathcal{F}(\mathbf{X};\mathbf{Y})\) 的一个子集 H（赋有一致收敛的一致结构），取 \(f\) 为典型映射 \((h,x)\to h(x)\) ；这时命题 2 的条件 2）自动满足，因为对 Y 的每个围域 W，由偶 \((h',h'')\) （满足 \((h'(x),h''(x))\in W\) 对一切 \(x\in X\) ）所成之集按定义是 H 的一致结构的一个围域。因此只需表述条件 1）；换句话说：

推论. 设 X、Y 是两个一致空间，H 是 \(\mathcal{F}(X;Y)\) 的一个子集。则 H 一致等度连续当且仅当映射 \((h,x) \to h(x)\) （由 \(H \times X\) 到 Y 中）一致连续，这里 H 赋有一致收敛的一致结构。

